% FILE: 01_research_question.tex
% PROJECT: otel-weaver-exp-001-custom-pi-registry
% THESIS ROLE: otel-semconv-evidence-surface-experiment

\section{Research Question}
\label{sec:otel-weaver-exp-001-custom-pi-registry-question}

\subsection{Problem Statement}
\label{sec:otel-weaver-exp-001-custom-pi-registry-question-problem}

Operational telemetry is not process evidence. This claim, stated in the doctrine file
\texttt{otel-weaver/doctrine/telemetry-is-not-process-evidence.md}, identifies a category error
that pervades observability-first software architectures: engineers instrument code with OpenTelemetry
spans, collect traces in a backend, and then treat the resulting data as a record of what the process
did. But a span that carries no Petri net token state, no witness hash, no lifecycle phase, and no
case identifier cannot be replayed as a lawful process trace. It is a measurement artifact, not a
process-evidence artifact.

The problem EXP-001 addresses is therefore definitional before it is technical: \emph{what is the
minimum schema a span must satisfy before it qualifies as admissible process feedstock?}
The experiment answers this question by constructing a machine-verifiable registry that the OTel
Weaver toolchain can validate, making the feedstock boundary explicit, externally auditable, and
version-pinned rather than implicit, undocumented, and drifting.

\subsection{Old-Regime Assumption Challenged}
\label{sec:otel-weaver-exp-001-custom-pi-registry-question-old-regime}

The old-regime assumption is that any span emitted by a correctly instrumented service constitutes
a record of what that service did at the process level. This assumption is embedded in the default
OTel data model, which makes attributes optional and imposes no structural constraints on how
spans relate to process lifecycle events, token states, or witness identities. Under this assumption,
conformance checking is impossible: there is no guarantee that two spans describing the same
activity carry comparable fields, and no mechanism to detect when a span has been emitted outside
a lawful object lifecycle.

The assumption also conflates schema registry drift with process drift.
The doctrine file \texttt{otel-weaver/doctrine/registry-diff-is-not-process-drift.md} (AUTHOR THESIS)
argues that a change in the OTel registry does not imply that the underlying process has changed;
the two must be tracked independently. EXP-001 challenges the old regime by making the registry
a typed, versioned artifact subject to Weaver validation, separating schema evolution from
process-conformance claims.

\subsection{Hypothesis}
\label{sec:otel-weaver-exp-001-custom-pi-registry-question-hypothesis}

A custom OTel Weaver semantic convention registry, binding eight named span attributes to required
or recommended cardinality under a version-pinned schema URL, can serve as the machine-verifiable
feedstock boundary for a process intelligence conformance court. Specifically:
\begin{enumerate}
  \item Spans that satisfy all four required attributes
        (\texttt{process.pi.instance\_id}, \texttt{process.pi.activity.name},
        \texttt{process.pi.lifecycle}, \texttt{process.pi.witness.hash},
        \texttt{process.pi.witness.id}) are admissible feedstock for the \texttt{BridgeRx}
        schema-bridging layer and the \texttt{validate\_and\_project} admission gate.
  \item Spans that violate one or more required attributes are refused at the gate and assigned
        a typed \texttt{RefusalCode} from the enum defined in \texttt{validator.rs}, preventing
        contaminated data from reaching the conformance court.
  \item The registry itself can be resolved and validated by the \texttt{weaver} CLI without
        modification to the standard OTel toolchain, demonstrating that process intelligence
        feedstock constraints are expressible in standard Weaver YAML.
\end{enumerate}

\subsection{Success Criteria}
\label{sec:otel-weaver-exp-001-custom-pi-registry-question-success}

Success requires all of the following:
\begin{itemize}
  \item \texttt{process\_pi.yaml} passes \texttt{weaver registry validate} without errors
        against a materialized \texttt{registry/} directory.
  \item The BLAKE3 receipt declared in \texttt{otel-weaver-source.manifest.yaml} is materialized
        to disk at \texttt{receipts/otel-weaver-source-receipt.json}, confirming the cryptographic
        chain is not merely declared but sealed.
  \item The \texttt{weaver\_integration\_tests.rs} test in \texttt{sources/wasm4pm/tests/}
        exercises \texttt{process\_pi.yaml} attributes directly, not only the Rust bridge/validator
        logic in isolation.
  \item An independent pm4py-based fitness/precision check is run against an event log produced
        from spans that satisfy the \texttt{process.pi.activity} schema, confirming that schema
        compliance at the feedstock layer propagates to conformance quality at the replay layer.
\end{itemize}
As of 2026-06-01, none of these criteria have been independently verified, placing the experiment
at PARTIAL status.
